\documentclass[10pt,a4paper]{amsart}
\usepackage[margin=19mm]{geometry}
\usepackage{amsmath,amssymb,mathtools}
\usepackage[scr=rsfs,cal=euler]{mathalfa}
\usepackage{xfrac}
\usepackage[unicode,hidelinks,bookmarks=false]{hyperref}
\newcommand{\ie}{i.e.\ }
\newcommand{\cf}{cf.\ }
\newcommand{\resp}{respectively\ }
\newcommand{\Subsection}{\subsection}
\providecommand{\nobreakdash}{\nobreak\hbox{-}}
\setlength{\emergencystretch}{2em}
\multlinegap0pt
\usepackage{amssymb}
\newtheorem{theorem}{Theorem}
\newtheorem{corollary}{Corollary}
\newtheorem{proposition}{Proposition}
\newtheorem{lemma}{Lemma}
\title{$\mathbb R^{\omega_1}$-Factorizable Spaces and Groups}
\author{Anton Lipin, Evgenii Reznichenko, Ol'ga Sipacheva}
\date{Source: arXiv:2509.05105v2 (CC BY 4.0)}
\begin{document}
\maketitle

\noindent\emph{Source and licence.} $\mathbb R^{\omega_1}$-Factorizable Spaces and Groups. Version arXiv:2509.05105v2 (CC BY 4.0), distributed by arXiv under the Creative Commons Attribution 4.0 International licence, http://creativecommons.org/licenses/by/4.0/, as recorded in the arXiv licence metadata for this exact version and shown on the abstract page https://arxiv.org/abs/2509.05105v2. Full licence notice: \texttt{licenses/arxiv-2509.05105-CC-BY-4.0.txt}.\par\smallskip

\noindent\textit{Editorial scope.} Counted unit: the article's proposition Proposition\_4.5 (source0.tex lines 811--815). The article's complete original proof of it is delivered with it (source0.tex lines 817--899, one \texttt{proof} environment of 83 lines). No mathematics is written, repaired or re-derived: the statement and the proof are the article's own lines, in the article's own notation, and no retained line is substituted or re-flowed.  Retained uncounted as the source-local context the proof needs: lines 492--501; lines 614--619.  Nothing was omitted from any retained span, so the package carries no omission record.  No retained line contains a citation, so the wrapper carries no bibliography and no bibliography entry.  No retained line contains a figure environment, an image-inclusion command or drawing code, so no image is delivered and the package declares no asset dependency.  Editorial formatting only, in addition: an AMS \texttt{amsart} preamble that restates the article's own declarations and macros used by the retained lines, namely the environments \texttt{theorem}, \texttt{corollary}, \texttt{proposition}, \texttt{lemma} on separate counters, together with the standard packages those lines need.  The retained environments print in this document as Theorem 1, Corollary 1, Proposition 1, Lemma 1 in order of appearance, whereas the article prints the same items with the numbering of its own full document.  The complete native source of the article as distributed in the arXiv e-print for this exact version is bundled unchanged as \texttt{sources/184536/source0.tex}.
\begin{theorem}
\label{Theorem_3.2}
Let $X$ be an $\mathbb R^{\omega_1}$-factorizable space. Then
\begin{enumerate}
\item[\rm(1)]
$X^\omega$ is hereditarily Lindel\"of and hereditarily separable; 
\item[\rm(2)]
if $w(X)\leq\omega_1$, then $X$ is second-countable.
\end{enumerate}
\end{theorem}

\begin{corollary}
\label{Corollary_4.1}
 A free filter $\mathscr F$  on a countable set $X$ is an $A$-filter if and only if 
the space $X_{\mathscr F}$ is $\mathbb 
R^{\omega_1}$-factorizable.
\end{corollary}

\begin{proposition}
\label{Proposition_4.5}
Any  $A$-filter of uncountable weight on $\omega$ contains an $A$-filter 
of weight $\omega_2$. 
\end{proposition}

\begin{proof}
Let $\mathscr F$ be an $A$-filter of uncountable weight on $\omega$. 
Then $w(\mathscr F)>\omega_1$ by Corollary~\ref{Corollary_4.1} and 
Theorem~\ref{Theorem_3.2}. 

\begin{lemma}
For any family $\mathscr A\subset \mathscr F$ with $|\mathscr A|\leq\omega_2$, there exists an 
$E(\mathscr A)\subset \mathscr F$ such that $\mathscr A\subset E(\mathscr A)$ and 
$w(E(\mathscr A))=| E(\mathscr A)|=\omega_2$. 
\end{lemma}
 
\begin{proof} 
If $w(\mathscr A)=\omega_2$, then we set $E(\mathscr A)=\mathscr A$. Suppose that $w(\mathscr 
A)<\omega_2$. We construct a family $\{M_\alpha:\alpha<\omega_2\}\subset \mathscr F$ completing $\mathscr A$ 
to the required family $E(\mathscr A)$ by induction on $\alpha<\omega_2$.

Take $\alpha<\omega_2$ and suppose that $\{M_\beta:\beta<\alpha\}\subset \mathscr F$ is already constructed. 
Let $\mathscr A_\alpha=\mathscr A\cup\{M_\beta:\beta<\alpha\}$. Since 
$w(\mathscr A_\alpha)\leq w(\mathscr A)+|\alpha|\le\omega_1$, it follows 
that $F(\mathscr A_\alpha)\neq \mathscr 
F$. Choose any $M_\alpha\in \mathscr F\setminus F(\mathscr A_\alpha)$. 

Having defined $M_\alpha$ for all $\alpha<\omega_2$, we set $\mathscr E=\mathscr 
A\cup\{M_\beta:\beta<\omega_2\}$. Note that $|\mathscr E|=\omega_2$. Let us show that $w(\mathscr 
E)=\omega_2$. Assume the contrary. Then $w(\mathscr E)<\omega_2$ and hence $F(\mathscr E)$ has a base $\mathscr 
B\subset F(\mathscr E)$ with $|\mathscr B|<\omega_2$. We have $\mathscr B\subset F(\mathscr A_\alpha)$ for 
some $\alpha<\omega_2$. Therefore, 
$$
F(\mathscr E)=F(\mathscr B)\subset F(\mathscr A_\alpha) \not\ni M_\alpha 
\in \mathscr E.
$$ 
This contradiction shows that $w(\mathscr E)=\omega_2$. We set $E(\mathscr A)=\mathscr E$. 
\end{proof}

\begin{lemma}
For any family $\mathscr A\subset \mathscr F$ with $|\mathscr A|\leq\omega_2$, there exists a $B(\mathscr 
A)\subset \mathscr F$ such that $\mathscr A\subset B(\mathscr A)$, $|B(\mathscr A)|\leq \omega_2$, and each 
$\mathscr C\subset \mathscr A$ with $|\mathscr C|\leq\omega_1$ has a countable base $\mathscr B\subset B(\mathscr 
A)$. 
\end{lemma}

\begin{proof}
Let us index (possibly with repetitions) the elements of $\mathscr A$ by ordinals: $\mathscr 
A=\{A_\alpha:\alpha<\omega_2\}$. For each $\alpha<\omega_2$, we set 
$\mathscr A_\alpha=\{A_\beta:\beta<\alpha\}$. 
Note that $w(\mathscr A_\alpha)\le\omega_1$. Since 
$\mathscr F$ is an $A$-filter, it follows that 
$\mathscr A_\alpha$ has a countable base $\mathscr B_\alpha\subset \mathscr F$ and that
the family $B(\mathscr A)=\mathscr A\cup
\bigcup\{\mathscr B_\alpha:\alpha<\omega_2\}$ is as 
required. 
\end{proof}

We proceed to prove the proposition. For each $\alpha<\omega_2$, we define 
a family $\mathscr A_\alpha\subset\mathscr F$ by induction as follows. We set $\mathscr A_0=\{\omega\}$. 
Assuming that $0<\alpha<\omega_2$ and $\mathscr A_\beta$ are already defined for all $\beta<\alpha$,  
we set 
$$
\mathscr A_\alpha^\ast=\bigcup_{\beta<\alpha}\mathscr A_\beta,\quad 
\mathscr E_\alpha=E(\mathscr A_\alpha^\ast),\quad \text{and}\quad 
\mathscr A_\alpha=B(I(\mathscr E_\alpha)).
$$

Having defined all $\mathscr A_\alpha$, we set 
$\mathscr A=\bigcup_{\alpha<\omega_2}\mathscr A_\alpha$ and $\mathscr 
G=F(\mathscr A)$. By construction, $|\mathscr A|=\omega_2$ and, therefore, $w(\mathscr G)\leq\omega_2$.

Let us show that $w(\mathscr G)=\omega_2$. Suppose that, on the contrary, $w(\mathscr G)<\omega_2$. 
Then $\mathscr G$ has a base $\mathscr B\subset \mathscr G$ with $|\mathscr B|\leq\omega_1$. Note that 
$\mathscr G=\bigcup_{\alpha<\omega_2}F(\mathscr A_\alpha)$. Therefore, $\mathscr B\subset 
F(\mathscr A_\alpha)\subset F(\mathscr E_{\alpha+1})\subset \mathscr G$ for some $\alpha<\omega_2$. 
Clearly, $\mathscr B$ is a base for $F(\mathscr E_{\alpha+1})$, which 
contradicts $w(\mathscr E_{\alpha+1})=\omega_2$.

It remains to show that $\mathscr G$ is an $A$-filter. Consider any $\mathscr M\subset \mathscr G$ with 
$|\mathscr M|\leq\omega_1$. From the same considerations as above, $\mathscr M\subset F(\mathscr E_\alpha)$ for some 
$\alpha<\omega_2$, and since $I(\mathscr E_\alpha)$ is a base for $F(\mathscr E_\alpha)$ and 
$|\mathscr M|\le \omega_1$, 
it follows that $\mathscr M$ has a base $\mathscr L\subset I(\mathscr E_\alpha)$ with 
$|\mathscr L|\leq\omega_1$. In turn, by the definition of $B(I(\mathscr E_\alpha))$, 
$\mathscr L$ has a countable base $\mathscr B\subset B(I(\mathscr E_\alpha))=\mathscr A_\alpha\subset \mathscr G$, 
as required. 
\end{proof}

\end{document}
